\section{The Lift \& Twist model}

The visible \(4\times4\) selector is not best understood as a scalar threshold.  It is a diagonal identity test seen through two different readout orders.

The C-side visible order is
\[
    O0,\ O1,\ B0,\ B1,
\]
where \(O\) denotes the ordinary shell, \(B\) denotes the branch shell, and the suffix denotes the shell rank.

The anchor-side visible order is
\[
    B0,\ O0,\ O1,\ B1.
\]

Thus the hidden identity rule
\[
    O0\mapsto O0,\quad
    O1\mapsto O1,\quad
    B0\mapsto B0,\quad
    B1\mapsto B1
\]
appears in visible coordinates as
\[
    0\mapsto 1,\quad
    1\mapsto 2,\quad
    2\mapsto 0,\quad
    3\mapsto 3.
\]

Equivalently, the visible permutation is
\[
    [1,2,0,3] = (0\ 1\ 2)(3).
\]

\begin{definition}[Lift \& Twist]
Lift \& Twist is the coordinate motion that lifts visible C-cycle and anchor-cycle data into native shell/rank coordinates, applies the identity lock
\[
    \mathrm{same\ shell} \wedge \mathrm{same\ rank},
\]
and returns the result to the visible register order, where it appears as a three-cycle plus one fixed branch state.
\end{definition}

In this reading, the twist is not decoration.  It is the cost of returning lifted shell/rank identity to the visible anchor ordering.
